% ================================================================
%  chapters/parts/ch01-c.tex  —  第 1 章 分片 c（本章末分片）
%  §1.2.2 3)（续）其他测孔容方法、§1.2.3 颗粒度测定、
%  §1.2.4 密度测定、§1.2.5 催化剂机械强度的测定
%  源：S1（上册），书内 p.61—80 / PDF p.80—99
%  接 ch01-b.tex（孔容和孔尺寸分布）之后
% ================================================================

\paragraph{补足上一节：两种不用吸附仪的孔容测法}除压汞法与气体吸附法之外，还有两种直接而
简便的方法可以测定催化剂的孔容。它们的准确性和精确度虽然不高，但在催化研究和催化剂制备
过程中有很强的实用性。

（1）\textbf{称重法}：把已知质量的干燥催化剂颗粒在蒸馏水中煮沸几分钟，或者置于真空瓶中
彻底抽空后浸满水，使水渗入孔内；再通过小心地过滤、气吹、沥离或者在湿布上碾滚等办法除去
颗粒间隙水，然后称重，由增重与水的已知密度计算孔容。

（2）\textbf{水滴定法}：将准确称量的干燥催化剂颗粒或粉末放入锥形瓶，用滴定管滴入蒸馏水，
不断摇晃并振荡锥形瓶使水与粉末混合均匀；当粉末突然成团或者粘在瓶底不动的瞬间即为吸水
饱和终点，由滴入水量计算孔容。

这两种方法测得的\textbf{吸水容量}可直接用于等体积浸渍法和初湿浸渍法制备催化剂时计算浸渍
液的配置浓度。这是它们真正的用途：在这里孔容不是物性表征量，而是配方计算量。也正因为
如此，它们对绝对准确度要求不高，却对「颗粒间隙水是否除净」极为敏感——间隙水未除净，就把
堆积空隙算成了孔容，结果系统性偏大；而水滴定法的终点判据是人为目视判断，重复性取决于
操作者。\src{S1 §1.2.2 3)，书内 p.61 / PDF p.80}

\subsection{颗粒度测定}

\begin{probout}
\textbf{本节回答}：「催化剂颗粒有多大」这句话在什么意义上可以被量化；一个粒径数值要连同
哪些条件才能被复核；以及为什么两台仪器测同一批样品可以差出两三倍。

\textbf{输出}：颗粒体系的\textbf{平均粒径}与\textbf{粒度分布}（含百分数基准、特征值
$D_{50}$、$D_{97}$、边界粒径与分布宽度指标），以及\textbf{所依据的等效粒径定义与测定
方法}。

\textbf{本节核心判断}：颗粒度测定测的不是几何尺寸，而是\textbf{等效量}。真实颗粒形状多
为不规则体，用一个或几个尺度数值去描述其尺寸特征是不现实的。因此表达颗粒度的方法放弃了
采用任何颗粒形状的概念，而引入\textbf{等效粒径}：当被测定颗粒的某种几何学特性（投影、
体积、表面积）、物理性质（光学、电学、质量）、物理行为（流体动力学、空气动力学）或者
化学吸附特性与某一直径的同质球体最相近时，就把该球体的直径作为被测颗粒的等效粒径。由此
推出本节两条结论：（i）\textbf{描述粒径及其分布必须同时说明依据的规则和测量的方法}，少了
任何一项，数值都不可复核；（ii）不同原理测的是\textbf{不同的等效直径}——体积等效、面积
等效、Stokes 等效、散射等效彼此不能换算，\textbf{换方法得到的数值不可直接互比}。
\end{probout}

\subsubsection{颗粒、一次颗粒与二次颗粒}

「颗粒」是指在毫米到纳米尺寸范围内具有特定形状的几何体；多相催化研究的一般是固体颗粒。
对单颗粒而言，颗粒度（particle size）指颗粒的大小。但在多相催化研究领域中，「颗粒」一词
所指的对象丰富而复杂。

狭义的催化剂（工业催化剂、实用催化剂）颗粒——球（sphere）、条（pellet、cylindrical）、
片、三叶草条（trilobe）等——是指人工成型的、具有发达孔系的颗粒聚集体。\textbf{工业催化剂
颗粒度}是指在操作条件下，催化剂颗粒\textbf{不再人为分开的最小基本单元}的大小或尺寸。
当催化剂颗粒具有显著的几何尺度（如毫米级）和确定的几何特征（如球体、圆柱体、环、三叶草
条等模状）时，催化剂的尺寸以其实际几何特征尺度标示——这是唯一可以直接用尺子量的情形。

更多的情况下，多相催化研究的「颗粒」以单晶、微粒、粉体及分散在载体上的金属或化合物粒子
等分散体系形态出现。其中：

\begin{itemize}
\item \textbf{一次颗粒}（primary particles）：原子、分子、离子按晶体结构的规则形成的单晶
粒子，通常是纳米级的；
\item \textbf{二次颗粒}（secondary particles）：若干晶粒聚集成的大小不一的微米级颗粒；
\item \textbf{粉体}（powder）：细小固体颗粒的集合。
\end{itemize}

这一区分在催化剂研究中是硬约束：X 射线衍射线宽化法测的是\textbf{晶粒}（一次颗粒），
显微镜法测的是\textbf{投影轮廓}，沉降法与激光散射法测的是\textbf{分散在介质中的颗粒或
团聚体}（二次颗粒）。三者测的常常不是同一个对象，因此三者的数值本来就不该相等
\src{S1 §1.2.3 1)，书内 p.61--62 / PDF p.80--81}。

\subsubsection{等效粒径：粒径必须连同方法一起报告}

\begin{principle}
\textbf{等效粒径的定义。}当被测定颗粒的某种几何学特性（投影、体积、表面积）、物理性质
（光学、电学、质量）、物理行为（流体动力学、空气动力学）或者化学吸附特性等性质与某一直径
的同质球体最相近时，就把该球体的直径作为被测颗粒的等效粒径。

\textbf{粒度分布。}研究的对象是一个粒子数量庞大的颗粒体系（集合），被研究的颗粒并不总是
绝对一样大小，有时大小分布甚至覆盖几个数量级。测量其中单颗颗粒的粒径没有意义，需要获得的
是关于颗粒体系\textbf{平均粒径}及\textbf{粒度分布}的总体信息。因此对颗粒体系而言，
「颗粒度」同时具有平均粒径大小和粒度分布双重含义。

\textbf{平均粒径的通式。}设粒度频率分布函数为 $f=q(x,n)$，$x$ 为粒径；$n=0$ 为个数分布
函数，$n=2$ 为面积分布函数，$n=3$ 为体积或质量分布函数，则
\begin{equation}\label{eq:ch01c-dn}
\bar{D}(n)=\frac{\int_{0}^{\infty}x^{n}f(x)\,\mathrm{d}x}{\int_{0}^{\infty}f(x)\,\mathrm{d}x}
\end{equation}
但更多情况下只能得到粒度\textbf{区间}分布的数据而非分布函数，这时平均粒径按下式计算：
\begin{equation}\label{eq:ch01c-dpq}
\bar{D}[p,q]=\left[\frac{\sum_{i}n_{i}d_{i}^{p}}{\sum_{i}n_{i}d_{i}^{q}}\right]^{\frac{1}{p-q}}
\end{equation}
式中 $n_i$ 为颗粒粒径分级第 $i$ 组级区间上的颗粒数；$d_i$ 为第 $i$ 组级颗粒的直径典型值，
取组级区间上、下限的几何平均值
\begin{equation}\label{eq:ch01c-di}
d_{i}=\sqrt{D_{i}D_{i+1}}
\end{equation}
$D_i$ 为该区间的下限粒径数值，$D_{i+1}$ 为上限粒径数值。

常用的几个特例是：个数平均粒径 $D[1,0]$（算术平均粒径）、个数表面积平均粒径
$D[2,0]$、个数体积平均粒径 $D[3,0]$、\textbf{表面积平均粒径} $D[3,2]$（Sauter 平均粒径）
与\textbf{体积（质量）平均粒径} $D[4,3]$（de Brouckere 平均粒径）。由
式~\eqref{eq:ch01c-dn} 可见它们与 $\bar{D}(n)$ 的关系是
\begin{equation}\label{eq:ch01c-d32a}
D[3,2]=\bar{D}(2),\qquad D[4,3]=\bar{D}(3)
\end{equation}
在很多情况下颗粒的粒径\textbf{个数}分布是得不到的，得到的是粒径面积、粒径体积或者粒径
质量分布，则此时
\begin{equation}\label{eq:ch01c-d32}
D[3,2]=\frac{\sum_{i}v_{i}}{\sum_{i}\left(v_{i}/d_{i}\right)},\qquad
D[4,3]=\frac{\sum_{i}w_{i}d_{i}}{\sum_{i}w_{i}}=\frac{\sum_{i}v_{i}d_{i}}{\sum_{i}v_{i}}
\end{equation}
式中 $v_i$（或质量 $w_i$）为颗粒粒径分级第 $i$ 组级区间上的总颗粒体积（质量），$d_i$ 仍
取该区间上、下限的几何平均值。

\textbf{这两式的意义。}$D[3,2]$ 是「总表面」与「总体积」之比决定的平均径，它是与比表面
积、传质、催化反应速率最直接相关的一个平均径；$D[4,3]$ 是体积加权平均径，它被大颗粒
支配，是激光粒度仪默认报出的那个数。两者相差多少，完全取决于分布的宽度：分布越宽，差别
越大。

\textbf{这两式的原书对应式。}此处已对照原页面图像（书内 p.63 / PDF p.82）：原书式 (1-54)、
(1-55) 印 $D[3,2]=\frac{\sum_{i}s_{i}d_{i}}{\sum_{i}s_{i}}$、
$D[4,3]=\frac{\sum_{i}w_{i}d_{i}}{\sum_{i}w_{i}}$，其变量说明段作「$s_{i}$ 和 $w_{i}$ 为颗粒粒径
分级第 $i$ 组级区间上的总颗粒面积和总颗粒体积（质量）」，可见原书确有面积形式的
$D[3,2]$。本处式~\eqref{eq:ch01c-d32} 改用体积（质量）分数写出，两者在全区间上等价：球形
颗粒的 $s_{i}\propto v_{i}/d_{i}$，形状系数在分子分母中相消，故
$\sum_{i}s_{i}d_{i}/\sum_{i}s_{i}=\sum_{i}v_{i}/\sum_{i}(v_{i}/d_{i})$；保留体积形式是为了能直接
对接激光粒度仪默认报出的体积基准分布。
\end{principle}

\begin{measure}
\textbf{必须同时报告的三件事。}（1）\textbf{等效粒径的定义}（体积等效、面积等效、
Stokes 等效、散射等效……）；（2）\textbf{分布所依据的百分数基准}——个数基准（count
basis）、面积基准（surface basis）、体积基准（volume basis）还是质量基准（mass
basis）；（3）\textbf{测定方法与制样分散条件}。同样一批 γ-氧化铝粉，图像法给出的是个数
基准的投影径分布，激光散射法给出的是体积基准的散射等效径分布，两者不仅\textbf{中位径
不同}，连\textbf{分布形状}都不同——因为从个数基准换算到体积基准要乘 $d^{3}$，这一步会
把大颗粒的权重急剧放大。\src{S1 §1.2.3 2)，书内 p.62--63 / PDF p.81--82}
\end{measure}

\subsubsection{粒度分布的表达与特征指标}

粒度分布（particle size distribution）就是用特定的方法表示颗粒体系中不同粒径颗粒占颗粒
总数（总量）的比例，反映粒子大小的均匀程度。表达上有两种基本形式：\textbf{频率分布}
（又称微分分布或区间分布，给出与各粒径相对应的颗粒在体系中所占的百分数）与\textbf{累积
分布}（又称积分分布，表示小于或大于某粒径的颗粒所占的百分数）。百分数的基准可用个数
基准、面积基准、体积基准和质量基准等表示——\textbf{同一组颗粒，换一个基准就是另一条
曲线}。

粒度分布可用表格、图形和函数三种形式表示：表格法即用列表方式给出某些粒径所对应的百分比；
图形法是在直角坐标系中用直方图和曲线等图形表示；函数法则用函数表示，如
Rosin--Rammler 分布函数、正态分布函数等。

表示粒度分布特性的关键指标有：

\begin{itemize}
\item \textbf{平均粒径}（mean particle size）：通过对粒度分布加权的平均粒度的量，即
式~\eqref{eq:ch01c-d32a} 与式~\eqref{eq:ch01c-d32} 给出的一族数值。
\item \textbf{$D_{50}$}：一个样品的累积粒度分布百分数达到 50\% 时所对应的粒径。它的物理
意义是粒径大于它的颗粒占 50\%，小于它的颗粒也占 50\%。$D_{50}$ 也称中位粒径或中值粒径
（median particle size）。\textbf{注意它对应的是「累积到一半」，不是「出现最多」}。
\item \textbf{$D_{97}$}：累计粒度分布数达到 97\% 时所对应的粒径，物理意义是粒径小于它的
颗粒占 97\%。$D_{97}$ 常用来表示粉体\textbf{粗端}的粒度指标；类似地还有 $D_{90}$、
$D_{3}$ 等。
\item \textbf{最频值}（mode）：频率曲线最高点所对应的粒径值，也称最可几粒径。
\textbf{最频值与 $D_{50}$ 只有在分布对称时才重合}；对常见的右偏分布（大颗粒拖尾），最频
值小于 $D_{50}$。
\item \textbf{边界粒径}：用来表示样品粒度分布的范围，由一对特征粒径组成，如
（$D_{10}$，$D_{90}$）或（$D_{3}$，$D_{97}$）。
\item \textbf{分布宽度}：工程上常报「跨度」（span）$=(D_{90}-D_{10})/D_{50}$。$D_{50}$
只说明「一半颗粒比它小」，跨度才说明「分布有多宽」。只报 $D_{50}$ 而不报跨度，等于只报了
分布的一个点。此处已对照原页面图像（书内 p.63 / PDF p.82）：原书「表示粒度分布特性的关键
指标」下确只列平均粒径、$D_{50}$、$D_{97}$、最频值（mode）与边界粒径五项——其中边界粒径举例
为「（$D_{10}$，$D_{90}$）、（$D_{3}$，$D_{97}$）」，$D_{97}$ 条附「类似地，还可以有 $D_{3}$、
$D_{10}$、$D_{90}$ 等」——原书未给出「跨度」；跨度 $=(D_{90}-D_{10})/D_{50}$ 是本节的工程惯用
补充，且原书对最频值只说「频率曲线的最高点所对应的粒径值，也称最可几粒径」，未涉及它与
$D_{50}$ 的关系。
\end{itemize}

\src{S1 §1.2.3 2)，书内 p.62--63 / PDF p.81--82}

\subsubsection{测定方法与各自的等效口径}

等效粒径的概念决定了颗粒度的测定已经失去了通常的几何学测量的意义，而转化为材料物理、
化学性能的测试分析。因此，描述粒径及其分布必须同时说明依据的规则和测量的方法。

实验室常用的颗粒度测定方法很多：筛分法、显微镜（图像）法、重力沉降法、离心沉降法、库尔特
（电阻）法、小角激光光散射法、动态光散射法、电镜法、氮气吸附（BET）法等。要选择合适的
测定方法，依次考虑三件事：首先，颗粒材料的物化性质；其次，其测量上下限要能够覆盖所关心
或者欲研究的颗粒尺寸范围；最后，测量条件尽可能接近材料的使用条件，即对材料的应用有充分
敏感性。\src{S1 §1.2.3 3)，书内 p.63--65 / PDF p.82--84}

表~\ref{tab:ch01c-methods} 把常用方法的等效口径、分布基准、适用范围与主要局限并列在一起。
读这张表请从右往左读：先确认样品的粒径落在哪一行的方法量程内，再看该行的「关键局限」是否
正好命中自己的样品，最后才看「等效粒径」一列——那一列才是数值能不能与文献比较的判据。

{\srcsize
\setlength{\tabcolsep}{3pt}
\begin{longtable}{@{}p{26mm}p{44mm}p{26mm}p{62mm}@{}}
\caption{实验室常用粒度分析方法及其特点。\textbf{读者应重点看第二列与第四列}：第二列说明
该法给出的是哪一种等效直径、以什么为基准的分布；第四列说明在什么情况下这个数值会失真。
凡是第二列不同或基准不同的两个数值，都不可直接互比。「适用粒径范围」一列的量程照录原书
表 1-10，实际可测范围还取决于仪器型号与制样条件。\src{据 S1 表 1-10 与 §1.2.3 3)，
书内 p.64—65 / PDF p.83—84 整理。}}
\label{tab:ch01c-methods}\\
\toprule
方法 & 等效粒径与分布基准 & 适用粒径范围 & 关键局限与适用条件 \\
\midrule
\endfirsthead
\toprule
方法 & 等效粒径与分布基准 & 适用粒径范围 & 关键局限与适用条件 \\
\midrule
\endhead
\bottomrule
\endlastfoot
筛分法 & 物理筛分等效径；\textbf{质量}（体积）分布 & 38 — 6000 µm（标准筛）；非标准筛可至 $3\times10^{5}$ µm & 下限 38 µm（Tyler 400 目），更细的颗粒测不了；颗粒能否通过取决于取向、筛面运动形式、物料载荷与筛分时间，重复性差；不适于黏性物质与易碎样品，杆状粒子误差大 \\
光学显微镜法 & 几何学等效径（投影面积圆等效径即 Heywood 径）；\textbf{个数}分布 & 约 1 — 6000 µm & 代表性差、重复性差；人工图像分析工作量大、随意性大、速度慢；形状规整且分布窄时测几百个颗粒即可，形状极不规则且分布宽时需测数万个 \\
电子显微镜法（SEM／TEM） & 几何学等效径；\textbf{个数}分布 & 约 0.001 — 5 µm（原表按倍率给出：光学显微镜 1 µm 起、SEM 0.01 µm 起、TEM 0.001 µm 起，均未给上限，此处上限为常见仪器的实用值） & 观察的颗粒数极少，代表性最差；TEM 测到的是晶粒（一次颗粒）与负载金属纳米粒子，不是二次颗粒；制样与取样误差常大于测量误差 \\
沉降法（重力） & Stokes 径（\textbf{流体动力学等效径}）；\textbf{质量}分布 & 10 — 300 µm & 只适用于球形颗粒的滞流沉降（$Re<1$，放宽也不超过 $2$）；仅适用于单分散体系，要求恒温；不规则颗粒的 Stokes 径小于几何学等效径；$d<0.5$~µm 时布朗运动严重干扰 \\
沉降法（离心／X 光透过） & Stokes 径；\textbf{质量}分布 & 0.1 — 300 µm & 离心沉降以离心加速度替代重力，拓宽了测量下限；可见光透过法下限 $>5$~µm（细颗粒散射），低能 X 射线可到 0.1 µm，但不适于高吸收 X 射线物质与强磁性材料 \\
电阻法（库尔特原理） & 球\textbf{体积}等效径；个数或体积分布 & 约 0.4 — 1200 µm & 单只小孔管只能测其孔径的 2\%—60\%，大小颗粒之比大于 30 时需多量程，归一化带来误差；大颗粒易堵小孔管；不适于多孔颗粒、密度大或憎水的样品；颗粒带电会使结果失真 \\
小角激光光散射法 & 等效散射径；\textbf{体积}分布（默认报 $D[4,3]$） & 0.02 — 3000 µm & 需预知颗粒折光指数（混合样品只能取主要物质的）；分辨率低，对较均匀的样品给不出精确分布；颗粒形状越不规则、分布越宽，误差越大；要求单散射，对浓度与分散有要求 \\
动态光散射法 & 流体动力学等效径；\textbf{强度加权}粒径分布（体积／数量分布须模型换算） & 0.0001 — 10 µm & 主要覆盖纳米与亚微米级；由 Einstein--Stokes 方程求径一步不引入与颗粒性质有关的参数（只需温度与黏度已知），但强度到体积／数量的换算须另设形状、光学性质与散射模型假设，不能与电镜个数分布直接比较峰面积；$z$ 平均径与 PDI 是累积量分析给出的平均量与宽度指标，与分布反演及其他平均径不可混用；要求样品池中颗粒浓度低并保持高分散；对入射光单色性要求高，须用激光光源 \\
电泳法 & 电泳迁移率 $\mu_{\mathrm{e}}$ 与 Zeta 电位 $\zeta$（\textbf{不给出粒径分布}） & 0.005 — 30 µm & 电泳光散射测的是迁移率，再由电泳模型求 $\zeta$；由 $\mu_{\mathrm{e}}$ 或 $\zeta$ 反演粒径必须附加含尺寸的模型并在报告中说明，薄双电层 Smoluchowski 极限 $\mu_{\mathrm{e}}=\varepsilon\zeta/\eta$ 与粒径无关（同 $\zeta$ 的 100 nm 与 200 nm 球给出同一 $\mu_{\mathrm{e}}$）；仪器同时报出 Zeta 与粒径时，粒径来自 DLS 通道或另行测定，不能把 $\mu_{\mathrm{e}}$ 当作唯一反演依据；只适用于能在液相中稳定分散且带电的颗粒；结果对介质离子强度、pH 与颗粒表面状态敏感，测量条件必须与使用条件一致 \\
BET 比表面积法 & \textbf{表面积}等效径 $D[3,2]$（单值） & 0.03 — 1 µm & 由 $D=6/(S\rho)$ 得到的只是一个平均径，\textbf{给不出任何分布}；富含微孔的颗粒必须用外比表面积计算，否则数值无意义 \\
空气透过法 & 表面积等效径（Kozeny--Carman 方程） & 0.2 — 75 µm & 相对测量方法，不能精确反映真实粒径；不适于多孔固体；一般只用于无孔大颗粒材料的工艺过程与产品质量监控（水泥、冶金、建材行业有标准） \\
X 射线衍射线宽化法 & Scherrer 等效径（\textbf{晶粒}体积平均尺寸） & 5 — 50 nm & 测的是晶粒（一次颗粒），不能反映二次颗粒的粒度；有微应变与仪器宽化时必须先分离 \\
化学吸附法 & 由分散度换算的活性组分粒径 & 约 0.1 — 10 nm & 利用化学吸附的选择性，只适用于负载型催化剂的活性组分；结果通常以\textbf{分散度}表示，无粒度分布数据 \\
\end{longtable}}

\noindent 上表已对照原页面图像逐格核对：原书表 1-10 为 9 行 4 列（方法／原理、等效粒径和
分布／工具或仪器／测试范围 µm），本表保留其中方法、等效粒径与分布基准、测试范围三列，
未保留「工具或仪器」列；「测试范围」一列的量程已按原表数值逐行改写（如电泳法
$0.005$—$30$ µm、动态光散射法 $0.0001$—$10$ µm、BET 比表面积法 $0.03$—$1$ µm、
筛分法 $38$—$6000$ µm 而标准筛以外的筛具可至 $3\times10^{5}$ µm）；表末的
X 射线衍射线宽化法与化学吸附法两行为本表所增，不在原表之内；「关键局限」列为正文各段
叙述的归纳，不是原表栏目。\src{S1 表 1-10，书内 p.64 / PDF p.83}

下面按原理把这条线索串起来。

\paragraph{筛分法}是一种最传统、应用最广的粒度测定方法，使颗粒通过不同尺寸的筛孔来测试
粒度。筛分法分干筛和湿筛，可以用单个筛子控制单一粒径颗粒的通过率（质量），也可以用一垛由
粗到细按筛号顺序叠起的筛子同时测量多个粒径颗粒的通过率，计算质量分数，获得\textbf{以质量
为基准}的筛分粒径分布及平均粒径。筛分法所得的\textbf{筛分径}（sieving diameter）以颗粒
通过和被截留的粗、细筛孔直径的均值表示，是典型的「筛分等效径」。

筛子有网筛和板筛两种，网筛的网目由金属丝按正方形编织而成。筛子的规格（筛号）用筛孔尺寸
表示；筛号还常用「目」（mesh）表示，「目」指在筛面的每英寸（1 英寸 = 2.54 厘米）长度上
开有的筛孔数目。由于编网所用筛丝直径不同，筛孔大小也不同，不同国家、不同行业的筛网规格有
不同标准，\textbf{同一目数筛子的筛孔尺寸可能不一样}。因此用目数标示筛子筛号时必须同时
标明筛孔尺寸。标准筛是按法令规格化的一套筛具，我国尚未颁布国家标准套筛，目前国际上常用
ISO 系列和美国 Tyler 系列；Tyler 系列最细的是 400 目，筛孔尺寸 38 µm。因此
\textbf{筛分法测定范围在 38 µm 以上}，粒度大于数毫米的物料所用的筛子一般为非标准筛。
筛分法的优点是工具便宜、操作简单；分析样品量大，代表性强，数据统计精确度高；能直接提供
粒径质量分布。缺点是下限只有 38 µm，不能分析细颗粒；难以分析黏性物质，不能分析易碎颗粒
样品；分析非规则形状（如杆状）粒子误差大；筛分结果受筛分时间影响大，重复性差；速度慢、
耗时。\src{S1 §1.2.3 3)1)，书内 p.65--66 / PDF p.84--85}

\paragraph{显微镜法}直接观察个别颗粒，根据颗粒投影像测得投影径，获得以个数、面积为基准的
粒度分布。光学显微镜可以测定微米级的颗粒，电子显微镜（SEM、TEM）可以测定纳米级的颗粒和
晶粒。投影径是几何学等效粒径，常用的表示方式是\textbf{投影面积圆等效粒径}
（equivalent circular diameter），即与粒子的投影面积相同圆的直径，也称为 Heywood 径。
其他几何学等效粒径还有：

\begin{itemize}
\item \textbf{Martin 径}：又称定方向等分径，是指在指定方向上，将颗粒投影面积等份分割的
弦线的长度；
\item \textbf{Feret 径}：任意方向上与颗粒投影面相切的一对平行线间的距离；通常指其中最宽
的一对平行线间的距离；
\item \textbf{Krumbein 径}：又称同方向最大径，指在指定方向上，颗粒投影轮廓内最长弦线的
长度。
\end{itemize}

对不规则颗粒，各粒径值之间一般存在关系：Feret 径 $\geq$ Krumbein 径 $\geq$ Heywood 径。
这一点有直接的实操含义：\textbf{同一张显微照片，用不同的「量法」就会得到系统性不同的
数值}，因此图像法报告必须写明所用的是哪一种投影径。为提高可靠性和代表性，显微镜法要求
观察和测量足够多的颗粒数目：形状规整、粒度分布窄的颗粒测几百个就足够；形状极不规则、
分布很宽的则需测数万个。显微镜法的优点是光学显微镜便宜，可直接观察颗粒形状和颗粒团聚，
透射电镜还可以测定负载型催化剂载体上金属或氧化物的纳米晶粒；缺点是代表性差，重复性差，
只适合测定粒度分布窄的样品，且操作烦琐、人工图像分析工作量大。\src{S1 §1.2.3 3)2)，
书内 p.66--67 / PDF p.85--86}

\paragraph{沉降法}建立在悬浮颗粒的动力学稳定性与粒径的关系上：颗粒较小时分子热运动产生的
布朗运动起主要作用，分散体系趋于稳定；颗粒较大时重力产生的沉降破坏分散体系的稳定性。
在自然重力状态下，若一个表面光滑的球形颗粒在静止流体中\textbf{滞流沉降}（颗粒所受重力、
浮力和运动阻力平衡，匀速沉降），则沉降速度与粒径的关系用 Stokes 公式描述：

\begin{equation}\label{eq:ch01c-stokes}
u=\frac{(\rho-\rho_{0})\,g\,d^{2}}{18\eta}
\end{equation}

式中 $d$ 为颗粒直径；$\eta$ 为流体黏度；$\rho$ 为颗粒真密度；$\rho_0$ 为流体密度；
$u$ 为颗粒的沉降速度；$g$ 为重力加速度。当雷诺准数 $Re=d u \rho_{0}/\eta<1$ 时沉降形态
为滞流（严格地应满足 $Re<0.2$，但也可放宽至 $Re<2$）。Stokes 公式表明\textbf{沉降速度与
颗粒直径的平方成正比}，大颗粒沉降快、小颗粒沉降慢；沉降法就是根据颗粒在液体中沉降速度
不同来测量粒径和粒度分布的方法。

Stokes 公式描述的是球形颗粒的滞流沉降行为。由于沉降阻力与颗粒受力的有效面积有关，对不规
则颗粒，其表面粗糙程度、形状及投影面积均影响沉降速度和沉降形态。由沉降法测得的粒径称为
\textbf{Stokes 径}，相当于在流体中与该颗粒具有\textbf{相同沉降速度}的同质球形颗粒的直径，
属于流体动力学等效粒径。对不规则颗粒体系，颗粒形状与球形的差异程度越大，沉降法测得的
Stokes 径与几何学等效粒径差异也就越大（Stokes 径偏小）。

一般对颗粒而言加速阶段时间很短，通常忽略，可以认为整个沉降过程是匀速的。沉降速度
$u=H/t$（$H$ 为沉降高度，$t$ 为沉降时间），由此自然重力状态下的 $d$--$t$ 函数为

\begin{equation}\label{eq:ch01c-dt-g}
d=\left[\frac{18\eta H}{(\rho-\rho_{0})\,g\,t}\right]^{1/2}
\end{equation}

在重力场中细颗粒的沉降速度很小，现代沉降仪大多引入\textbf{离心沉降}方式：离心加速度远
大于重力加速度，粒径相同时离心沉降的测试时间大大缩短。离心沉降速度的表达式在形式上与
重力沉降式一样，只是以离心加速度代替重力加速度：

\begin{equation}\label{eq:ch01c-centrif}
u_{r}=\frac{(\rho-\rho_{0})\,\omega^{2}r\,d^{2}}{18\eta}
\end{equation}

式中 $\omega$ 为转动的角速度，$r$ 为颗粒距转动中心的距离，$u_r$ 为颗粒的离心沉降速度。
离心沉降时颗粒的实际运动轨迹是一个半径逐渐扩大的螺旋线，但由于颗粒和流体同时做圆周
运动，二者在切向上相对静止，所以\textbf{离心沉降速度并不是颗粒的实际运动速度，只是其在
径向上的分量}。当颗粒匀速沉降时离心力与摩擦阻力达到平衡；经时间间隔 $t$，颗粒由距转动
中心 $r_1$ 运动至 $r_2$ 处，积分整理后离心力状态下的 $d$--$t$ 函数为

\begin{equation}\label{eq:ch01c-dt-c}
d=\left[\frac{18\eta\,\ln\left(r_{2}/r_{1}\right)}{(\rho-\rho_{0})\,\omega^{2}t}\right]^{1/2}
\end{equation}

直接测量颗粒的沉降速度很困难，但通过测量颗粒体系沉降过程的\textbf{浓度变化}可以间接求得
沉降速度。依据浓度测定方法，常见的沉降测定法有移液管法、沉降天平法、沉降管法、比重计法、
光透过重力（离心力）沉降法和 X 射线透过重力（离心力）沉降法等。其中移液管法
（pipette analysis）是把颗粒均匀分散的沉降液转入带固定位置移液管的 Andreassen 容器，令其
自由沉降，在适当时间从一定高度处抽吸一定量的悬浊液，烘干并精确称量其中的固体质量，按
Stokes 公式把时间换算成粒径，由干重变化求粒度分布；光透过法则令沉降液自由沉降，用光强
变化率间接反映沉降速度，光的透过量与浓度的关系符合 Lambert--Beer 定律。
\src{S1 §1.2.3 3)3)，书内 p.67--70 / PDF p.86--89}

\begin{measure}
\textbf{沉降法的条件清单。}（1）\textbf{浓度}：沉降液中颗粒质量分数最好控制在 0.1\%—3\%，
否则颗粒间的相互作用会干扰沉降过程；颗粒间距离大于粒径 10 倍、沉降容器尺寸大于粒径
100 倍时，颗粒干扰和器壁效应可忽略。（2）\textbf{介质匹配}：被测颗粒不应在分散介质中
溶解、溶胀；需调节分散介质的密度和黏度与固体颗粒的密度和尺寸相匹配；不易润湿的颗粒还需
加入表面活性剂增大固液亲和力。（3）\textbf{滞流条件}：对于密度或粒径很大的颗粒，沉降
速度太快会进入过渡流或湍流区，Stokes 公式不再适用，体系\textbf{最大}颗粒的沉降雷诺准数
必须满足滞流条件。（4）\textbf{温度}：测定过程要求维持恒定温度。（5）\textbf{下限}：
非常微细的颗粒（如 $d<0.5$~µm）在介质中的布朗运动会严重干扰沉降，沉降法也无法测定；
重力沉降的粒径范围为 10—300 µm，离心沉降为 0.1—10 µm。光透过法的悬浮液质量百分浓度
以 0.02\%—1\% 为好，可见光检测细颗粒时发生光散射，测量下限大于 5 µm，且可见光强度较弱、
灵敏度低；选用低能 X 射线作入射光源可把下限推到 0.1 µm，但不适于分析高吸收 X 射线物质
和强磁性材料。\src{S1 §1.2.3 3)3)，书内 p.69--70 / PDF p.88--89}
\end{measure}

\paragraph{电阻法}的工作原理是：悬浮在电解液中的不导电颗粒随电解液流入一个负压的孔管，
当其流经小孔时，颗粒体积排除了相同体积的电解液，导致孔口处电阻发生瞬时变化；在孔管内外
设有电极，通过一个恒电流电路记录产生的电位脉冲。库尔特原理指出\textbf{电阻变化的大小与
颗粒的体积成正比}，因此电位脉冲信号的大小和次数与颗粒的体积大小和数目成正比。库尔特计数
法获得的是\textbf{球体积等效粒径}，其分布以个数或体积为基准。

库尔特计数法单只小孔管（对应一个量程）可测量的颗粒尺寸范围在其孔直径的 2\%—60\%；当
样品粒度分布较宽（大、小颗粒之比大于 30）时需使用多只小孔管分别测量。仪器提供的孔直径
一般为 15—2000 µm，量程约 0.4—1200 µm，测量速度最高可达每秒 10000 个颗粒。它以体积
等效粒径表征颗粒粒径，\textbf{不受颗粒形状影响}；它一个个地测量每个颗粒的大小，因而
分辨率很高，并能给出颗粒的绝对数目，特别适用于检测液体中的稀少颗粒；重复性好，分析不受
分散剂、悬浮液折光率、颜色、形态的影响。缺点是对样品要求苛刻：必须可分散并悬浮在电解液
中，因此无法测量相对密度较大的物质颗粒和憎水的有机材料；如果样品颗粒带电，会抵消其体积
占用造成的电阻增加，造成测量失真；动态范围小，宽分布样品虽可多管测量，但数据归一化会
带来误差，且大颗粒容易堵塞小孔管（孔口越小越易堵）；\textbf{不适宜测量多孔性颗粒}。
\src{S1 §1.2.3 3)4)，书内 p.70--71 / PDF p.89--90}

\paragraph{小角激光光散射法}（又称激光光衍射法）属于静态（弹性）光散射，仪器常称为激光
粒度仪。光束行进方向上的颗粒障碍物对光有散射作用，会使一部分光偏离原来的传播方向而向
四周传播；散射角的大小与颗粒尺寸有关——\textbf{颗粒越小，散射角越大}，而散射光的强度
包含颗粒的含量信息。该法把颗粒均匀地放置到激光光束中，用光电探测器在不同角度上测量散射
光强度，数据经理论处理得到粒度分布。

散射理论分两支。对于粒径远大于入射光波长的球形颗粒体系，可以认为所有颗粒对光线具有相同
的散射效率，颗粒的透光性也可忽略，此时 \textbf{Fraunhofer 衍射理论}是解析光散射数据的
简便方法：大颗粒对光的散射角很小，散射光强度与颗粒投影面积（粒径的平方）和数量成正比。
当颗粒直径接近入射光波长时，更多的光被散射到很大的角度，散射光强度与颗粒大小、折光率及
入射光强度之间存在复杂关系；而当颗粒直径远小于入射光波长时（Rayleigh 区），散射变成各向
同性，此时在材料、介质、波长与颗粒个数都相同的条件下，散射光强度与粒径的六次方\textbf{成
正比}——粒径由 10 nm 增至 20 nm，散射强度增大 $2^{6}=64$ 倍。原书此处排印为「成反比」，
与 Rayleigh 散射的 $I\propto d^{6}$ 关系相反，属原书笔误，本版按正比表述并在此记明。这时需用
\textbf{Mie 散射理论}描述散射光场强度分布——它由
Maxwell 电磁波方程严格推导，适用于计算任意大小的颗粒，可在很宽的粒度范围内得到完全精确的
结果（约 0.02—3000 µm），而 Fraunhofer 理论只是粒径远大于光波长情况下散射现象的近似和
简化描述。早期受计算能力限制，激光光衍射分析法只能使用 Fraunhofer 近似（这也是该法称为
「衍射法」的原因），探测器阵列角度也较小，只能测微米以上颗粒；目前先进仪器大多全量程使用
Mie 理论，同时采用多激光光源（短波长）、增加大角与背散射光电探测器、偏振光强度差散射
（PIDS）等技术改善分辨率与下限。

Mie 散射理论假设散射颗粒是球形，散射光强度与颗粒体积相关；小角激光光散射法正是利用颗粒
对激光的散射特性做等效对比来测定\textbf{等效散射径}（与实际被测颗粒具有相同散射效果的
球形颗粒的直径）。因此实际测量中不规则颗粒的形状、粒径分布特性对最终结果影响较大，颗粒
形状越不规则、粒径分布越宽，分析结果误差越大（引入非球形颗粒校正的 Mie 公式可用于不规则
颗粒）。该法得到的是\textbf{体积分布}，平均粒径是体积平均粒径 $D[4,3]$。由于 Mie 理论
假设散射光是不相关的单散射，实际测量对样品池中颗粒的浓度和分散情况有一定要求，以降低复
散射或相干散射的影响；小颗粒的散射与其材料的光学性质关系密切，Mie 理论需要预知颗粒的
折光指数（有系数匹配、折光仪和平均值等确定方法，混合样品可取主要物质的折光指数）。该法
测量速度快、动态范围大、操作简便、重复性好、应用领域广（不仅能测固体颗粒、乳胶颗粒，而且
可测雾滴、气泡等的粒度），但\textbf{分辨率低}，对于颗粒较均匀的样品不能给出精确的粒径
分布。\src{S1 §1.2.3 3)5)，书内 p.71--73 / PDF p.90--92}

\paragraph{动态光散射法}（dynamic light scattering，也称准弹性光散射 QELS 和光子相关
光谱 PCS）是研究纳米颗粒粒度分布的有力手段（1—500 nm）。前述小角激光光散射法可以理解为
静态光散射；对于远小于光波波长的亚微米级颗粒（$<250$~nm），由瑞利散射理论可知，散射光
相对强度的角分布与粒子大小\textbf{无关}，不能通过对散射光强度的空间分布来确定颗粒粒度，
而动态光散射正好弥补这一不足。

纳米颗粒在流体中布朗运动的速度与颗粒的大小有关，符合 Einstein--Stokes 方程
\begin{equation}\label{eq:ch01c-einstein}
D_{\mathrm{d}}=\frac{kT}{6\pi\eta R}
\end{equation}
式中 $D_{\mathrm{d}}$ 为扩散系数；$k$ 为玻尔兹曼常量；$T$ 为热力学温度；$\eta$ 为流体的
动力学黏度系数；$R$ 为颗粒的水力学半径。由于多普勒效应，散射光会产生微小的频率漂移
（多普勒频移）而谱线增宽，散射光频谱半分宽 $\Gamma$ 与扩散系数之间存在如下关系
\begin{equation}\label{eq:ch01c-doppler}
\Gamma=\left(\frac{4\pi n}{\lambda}\right)^{2}\sin^{2}\left(\frac{\theta}{2}\right)D_{\mathrm{d}}
\end{equation}
式中 $n$ 为流体的折光指数，$\lambda$ 为入射波长，$\theta$ 为散射角。光谱强度涨落和频率
漂移是互相关联的，研究这种由于颗粒运动而产生的散射光的强度涨落和精细结构，称为动态光
散射。通过相关技术精确测量颗粒散射光强与时间的函数关系并进行相关运算，可以得出颗粒粒径
及粒度分布。

由于动态光散射研究的是粒子的动态行为，它得到的也是一种\textbf{流体动力学等效径}。这里要
分清两件事：由相关函数衰减得到扩散系数 $D_{\mathrm{d}}$、再由 Einstein--Stokes 方程求流
体动力学径这一步，确实不引入与颗粒性质有关的参数（只需已知温度与介质黏度）；但仪器默认
报出的是\textbf{强度加权}的粒径分布，\textbf{体积}分布与\textbf{数量}分布都是在此基础上的
模型换算，须另设颗粒形状、光学性质与散射模型假设，不能与电镜的个数分布直接比较峰面积。
两者不可混用的算例：同种材料、Rayleigh 条件下等个数混合的 10 nm 与 20 nm 颗粒，体积权重之
比为 $1:8$，强度权重之比为 $1:64$（大颗粒分别占 $8/9$ 与 $64/65$）。此外，累积量分析给出
的 $z$ 平均径与多分散指数 PDI 是\textbf{平均量与宽度指标}，不是分布本身，也不等同于
$D[4,3]$、$D[3,2]$ 等平均径，报告中必须注明口径。该法不考虑散射光的复散射和纳米颗粒之间
的相互作用，因此要求样品池中颗粒的浓度较低并保持高分散，同时对入射光的单色性要求很高，
因而也以激光为光源。\src{S1 §1.2.3 3)6)，书内 p.73--74 / PDF p.92--93}

\paragraph{其他方法。}\textbf{电泳法}使用 Zeta 电位分析仪，在电场力作用下带电颗粒在悬浮
体系中定向迁移。它测的是\textbf{电泳迁移率} $\mu_{\mathrm{e}}$，再由电动力学模型得到
\textbf{Zeta 电位} $\zeta$，\textbf{并不直接给出粒径}——注意与动态光散射分开：后者测的是
布朗运动引起的扩散系数，由 Einstein--Stokes 方程求流体动力学径。把 $\mu_{\mathrm{e}}$ 或
$\zeta$ 反演成粒径，必须另加一个\textbf{含粒径的模型}，并在报告中写明该模型及适用条件。
反例：薄双电层 Smoluchowski 极限 $\mu_{\mathrm{e}}=\varepsilon\zeta/\eta$ 与粒径无关，
$\zeta=25$ mV 时 100 nm 与 200 nm 的球给出相同的 $\mu_{\mathrm{e}}$，可见迁移率本身不含
粒径信息。因此，当仪器同时报出 Zeta 电位与粒径时，粒径来自 DLS 通道或另行测定，
\textbf{不能}把电泳迁移率当作唯一的粒径反演依据。可测 0.005—30 µm 的颗粒。\textbf{BET 比表面积
法}通过比表面积求得与欲测粒子具有等比表面积的球形颗粒的等效粒径
\begin{equation}\label{eq:ch01c-betd}
D=\frac{6}{S\rho}
\end{equation}
其中 $S$ 为比表面积，$\rho$ 为颗粒的密度，形状因子对球体而言为 6。该法只能给出颗粒的
\textbf{表面积平均粒径} $D[3,2]$，而无法给出粒度分布；显然对富含微孔的颗粒应该用其
\textbf{外}比表面积计算平均粒径（否则得到的是孔壁面积对应的假想粒径）。
\textbf{空气透过法}通过测量颗粒床层的压降、改变床层高度使床层空隙率发生变化，用半经验的
Kozeny--Carman 方程计算颗粒的比表面积及平均粒径，但不能测量粒度分布；该法是一种相对测量
方法，不能精确反映颗粒的真实粒径，且不适于测定多孔固体，一般只用于无孔大颗粒材料生产的
工艺过程和产品质量监控（水泥、冶金、建材行业有相应标准）。
\textbf{X 射线衍射线宽化法}利用小晶粒的衍射峰宽化效应，根据 Scherrer 公式计算 5—50 nm
晶粒的体积平均大小（晶粒度），\textbf{并不能反映通常意义上的「颗粒」（二次颗粒）的粒
度}。\textbf{化学吸附法}利用化学吸附的选择性测定负载型催化剂表面活性组分的分散度和表面
积，进而根据一定的颗粒形状模型计算活性组分的平均粒径：粒径 $d$（单位 nm）与其分散度 $D$
之间可以简单地换算，$d\approx0.9/D$。\src{S1 §1.2.3 3)7)，书内 p.74--75 / PDF p.93--94}
此处已对照原页面图像（书内 p.75 / PDF p.94）：原书印为「对于负载型催化剂表面的金属颗粒，
粒径 $d$（单位 nm）与其分散度 $D$ 之间可以简单地换算：$d\approx0.9/D$」，常数 $0.9$ 为
原书原数，本节照录（OCR 把此处误识为「ｄ ．．．． ９／Ｄ」）。按原书式，$D$ 应以
\textbf{分数值}代入（$D=1$ 即全部原子暴露时 $d\approx0.9$ nm），不是百分数，故不必再引入
可调常数 $k$。

\subsubsection{同一批颗粒的三条曲线：一个读图实例}

\begin{figure}[htbp]
\centering
\begin{tikzpicture}
\begin{axis}[
  width=\onecolwidth, height=52mm,
  xmode=log, log basis x=10,
  xmin=0.3, xmax=40, ymin=0, ymax=1.08,
  domain=0.3:40, samples=400,
  xtick={0.5,1,2,5,10,20,40}, xticklabels={0.5,1,2,5,10,20,40},
  ytick={0,0.2,0.4,0.6,0.8,1.0},
  xlabel={等效直径 $D$ / µm},
  ylabel={归一化频率},
  tick label style={font=\fontsize{8}{9.5}\selectfont},
  label style={font=\fontsize{8.5}{10}\selectfont}]
% 个数基准：对数正态，ln D 的众数为 ln 2.2、几何标准差 exp(0.55)
\addplot[domain=0.3:40,samples=400,mainline,black]
  {exp(-((ln(x)-0.788)^2)/(2*(0.55^2)))};
% 体积（质量）基准：同批颗粒按 D^3 加权，峰移到 exp(0.788+3*0.55^2)=5.45 µm
\addplot[domain=0.3:40,samples=400,thinline,black,dashed]
  {exp(-((ln(x)-0.788)^2)/(2*(0.55^2)))*(x^3)/41.5};
% 沉降法：同一批颗粒，Stokes 径 d_St = 0.78 d_V（形状因子 0.78）
\addplot[domain=0.3:40,samples=400,thinline,black,densely dotted]
  {exp(-((ln(x/0.78)-0.788)^2)/(2*(0.55^2)))*((x/0.78)^3)/41.5};
\draw[accentsoft,dashed,thinline] (axis cs:2.2,0) -- (axis cs:2.2,1.08);
\draw[accentsoft,dashdotted,thinline] (axis cs:4.25,0) -- (axis cs:4.25,1.08);
\draw[accentsoft,densely dotted,thinline] (axis cs:5.45,0) -- (axis cs:5.45,1.08);
\node[anchor=north west,font=\fontsize{8}{9.6}\selectfont,align=left,
      text width=36mm] at (axis cs:0.32,0.98)
  {实线 图像法（个数基准）\\虚线 体积基准（同批换算）\\点线 沉降法 Stokes 径\\中位径 2.2 / 5.4 / 4.2 µm};
\end{axis}
\end{tikzpicture}
\caption{同一批颗粒、三种等效口径给出的三条分布曲线。读者应看三处：\textbf{峰的位置}
（个数基准峰在约 $\qty{2.2}{µm}$，同一批颗粒换算到体积基准后峰移到约 $\qty{5.4}{µm}$；
沉降法用 Stokes 径作横轴时峰落在约 $\qty{4.2}{µm}$）；\textbf{形状与宽度的变化}（体积基准
曲线相对个数基准明显右移并变窄，这是 $d^{3}$ 加权把大颗粒的权重放大所致，不是样品变了）；
\textbf{三条竖直参考线}（分别对应三种口径下的中位径，跨度为 2.2—5.4 µm，相差约 2.5 倍）。
三条曲线各自归一到峰值 1，因此\textbf{只能比较峰的位置，不能比较峰的绝对高度}。本图为按
对数正态个数分布及其 $d^{3}$ 加权换算构造的模型曲线（横轴为等效直径 $D$，形状因子取
0.78），用来说明「同一批颗粒换一种等效定义就换一个数」，不代表任何具体样品的实测数据。
\modelmark
\figsrc{据 S1（上册）§1.2.3 2)3)，书内 p.62--64 / PDF p.81--83 所述等效粒径与分布基准
的差异计算绘制。}}
\label{fig:ch01c-equivalent}
\end{figure}

\begin{obs}
在图~\ref{fig:ch01c-equivalent} 上逐项读数：实线（图像法个数基准）的峰在横轴
$\qty{2.2}{µm}$ 处，峰高 1；虚线（同一批颗粒按体积加权）的峰在约 $\qty{5.4}{µm}$ 处，
峰高 1；点线（沉降法，横轴为 Stokes 径）的峰在约 $\qty{4.2}{µm}$ 处，峰高 1。三条竖直
参考线分别落在 $\qty{2.2}{µm}$、$\qty{4.25}{µm}$、$\qty{5.45}{µm}$。虚线与点线的形状
相同，只是整体沿横轴平移了 $\lg(1/0.78)\approx0.11$ 个数量级；实线则明显更宽、更靠左。
三者在横轴 0.3—1 µm 区间都趋近于零；虚线与点线在 15 µm 以上也已趋近于零。
\end{obs}

\begin{infer}
由上述读数只能推出下面四条，其余不能推。

（1）\textbf{「同一批颗粒」的中位径可以是 2.2 µm，也可以是 5.4 µm}——差别不来自样品，
而来自「按什么加权」。若报告只写「平均粒径 2.2 µm」而不写基准，读者无法判断这是个数中位
径还是体积中位径，两个数相差 2.5 倍。

（2）虚线与点线形状完全一致、只差一个尺度因子 0.78，说明这一处的差别\textbf{不是分布
形状的差别，而是横轴定义（等效直径口径）的差别}：同一批颗粒，用沉降法量出的「直径」系统
性地比体积等效径小。这正是「Stokes 径与几何学等效粒径不可互比」的图形表达。

（3）三条竖线给出的是三个可比性完全不同的数：2.2 µm（个数基准）、5.4 µm（体积基准）、
4.2 µm（Stokes 径）。若把它们并列写进同一张表当作「同一批样品在三种仪器上的粒径」，读者
无从区分是样品不同还是方法不同——这是粒度数据最常被误用的一种方式。

（4）本图不能用来推断任何真实样品：0.78 这个形状因子是模型设定值，不是普适换算系数；
曲线之所以光滑，是因为它由解析函数生成。真实样品的差别还叠加了团聚、形状分布与仪器
响应函数的影响。\modelmark
\end{infer}

\begin{dataflow}
\textbf{从仪器信号到可报告的粒径数。}（1）\textbf{原始量}：筛分法得到各级筛上、筛下的
\textbf{质量}；沉降法得到不同时刻的\textbf{浓度或光强}；电阻法得到\textbf{电位脉冲}的
幅度与计数；激光散射法得到各角度的\textbf{散射光强}；图像法得到每个颗粒的\textbf{投影
轮廓}。（2）\textbf{换算}：用各自的理论把原始量转为某一等效直径的粒径分级数据——沉降法
用式~\eqref{eq:ch01c-dt-g} 或式~\eqref{eq:ch01c-dt-c} 把时间换算成 Stokes 径，电阻法用
脉冲幅度换算成体积等效径，激光散射用 Fraunhofer 或 Mie 理论反演成散射等效径。
（3）\textbf{加权}：得到的是某一\textbf{基准}的分布（筛分与沉降是质量基准，图像与电阻
可给个数基准，激光散射给体积基准）。（4）\textbf{求特征值}：用
式~\eqref{eq:ch01c-dpq} 计算所需平均径，或直接读取 $D_{50}$、$D_{10}$、$D_{90}$、$D_{97}$
与最频值。（5）\textbf{报告}：必须连同等效粒径定义、分布基准、方法名与主要条件一起给出。
任何一步跳过，第（5）步的数字就不可复核。\src{S1 §1.2.3 2)3)，书内 p.62--64 /
PDF p.81--83}
\end{dataflow}

\begin{pitfall}
\textbf{误判一：把不同方法（或不同基准）的粒径直接比较。}图像法个数基准
$\qty{2.2}{µm}$ 与激光散射体积基准 $\qty{5.4}{µm}$ 是同一批样品的两个不同等效量，比大小
没有意义；「换了一台仪器，粒径涨了一倍」通常只是基准或口径变了。判据：两个数值要能比较，
必须\textbf{等效粒径定义与百分数基准都相同}，且方法原理在样品的形状与分布宽度下成立。

\textbf{误判二：用单一平均径代替分布。}$D_{50}$ 只说「一半颗粒比它小」。宽分布样品中，
$D[3,2]$ 与 $D[4,3]$ 可以相差数倍，两者对应的是完全不同的物理量（前者关联比表面积与
传质，后者被大颗粒支配）。只报一个「平均粒径」而不报跨度，无法判断样品是否均匀。

\textbf{误判三：忽略量程与假设的越界使用。}筛分法下限 38 µm，报出更细的筛分径必然是外推；
沉降法在 $Re>2$ 或 $d<0.5$~µm 时给出的 Stokes 径不成立；库尔特法不宜测多孔颗粒；X 射线
线宽化给出的是晶粒（一次颗粒）尺寸，把它当作「催化剂颗粒度」是最常见的张冠李戴。

\textbf{误判四：取样与制样不代表整体。}颗粒样品粒子数量庞大而且很不均匀，实践表明，
粒度分析时\textbf{由采样和制样分散步骤带来的误差往往比分析的测定误差更大}。大颗粒体系的
主要误差来源是采样不均匀；纳米级细微颗粒的关键是分散（团聚体会被当成一个大颗粒，从而使
分布整体右移）。所有粒度分析方法都要求颗粒之间不能相互干扰。工业催化剂床层的上段与下段
经历过的磨损、破碎与积炭不同，取样位置必须写明。
\src{S1 §1.2.3 3)，书内 p.65、p.67 / PDF p.84、p.86}
\end{pitfall}

\begin{crosscheck}
\textbf{三条互补回查。}（1）\textbf{量级互查}：由氮吸附比表面积算出的表面积等效径
$D=6/(S\rho)$ 与图像法、沉降法的数值只能比较\textbf{数量级}，不能比较具体数值——三者的
等效定义不同；若相差超过一个数量级，应先怀疑某一个测错了（如 BET 面积含微孔、图像法把
团聚体当成单个颗粒）。（2）\textbf{分布互查}：同一样品用体积基准与个数基准各测一次，
两组数据之间应该满足 $d^{3}$ 加权的换算关系；若换算后峰位对不上，说明其中一次测量有团聚
或分割（sieve 堵塞、孔管堵塞）造成的假峰。（3）\textbf{与制备记录回查}：催化剂成型前后的
粒径变化应与成型工艺一致；二次颗粒的粒径若与一次颗粒（XRD 晶粒尺寸）相差不到两倍，通常
说明样品已经严重团聚或 XRD 峰宽主要来自微应变而非尺寸。
\src{S1 §1.2.3 2)3)7)，书内 p.62--64、p.75 / PDF p.81--83、p.94}
\end{crosscheck}

\subsection{密度的三种定义与测定}

\begin{probout}
\textbf{本节回答}：同一批多孔催化剂为什么会有三个（甚至更多）都叫「密度」的数；每一个数
究竟用哪一份体积做分母；它们各自怎么测；以及如何用它们算出孔隙率与床层空隙率。

\textbf{输出}：\textbf{骨架密度}（真密度）、\textbf{颗粒密度}、\textbf{堆积密度}三个数，
以及由它们派生的\textbf{孔隙率} $p$、\textbf{空隙率} $\varepsilon$ 与\textbf{比孔容}
$v_{\mathrm{p}}$。

\textbf{本节核心判断}：密度是物质的质量与体积之比，但对多孔固体催化剂，与其表面积一样，
其所占有的空间（体积）也因度量方法的不同而有不同的含义和量值。多孔催化剂颗粒堆积时的
\textbf{外观体积}实际包括三部分：颗粒本身的三维骨架、颗粒内部孔道、堆积时颗粒之间的空隙。
因此\textbf{以不同的体积除质量，所得密度的概念也就不同}。使用任何一个密度数值之前，必须
先问三件事：这个数\textbf{测的是哪一种}体积、用什么方法测的、样品处于什么堆积状态。文献
习惯及不同的组织和机构对固体催化剂各种密度的命名或定义并不一致，本节遵循文献主要传统并
参照 ASTM 的有关定义（ASTM D3766—08（2013）等）。
\end{probout}

\subsubsection{三种体积，三种密度}

\begin{principle}
\textbf{体积的分解。}催化剂颗粒堆积时的外观体积包括骨架、内孔与颗粒间空隙：
\begin{equation}\label{eq:ch01c-volapp}
V_{\mathrm{app}}=V_{\mathrm{void}}+V_{\mathrm{pore}}+V_{\mathrm{skel}}
\end{equation}
式中 $V_{\mathrm{app}}$ 为催化剂外观体积（堆积体积）；$V_{\mathrm{void}}$ 为堆积时颗粒之间
的空隙体积；$V_{\mathrm{pore}}$ 为颗粒内部\textbf{开孔}所占的体积；$V_{\mathrm{skel}}$ 为
固体骨架体积，包括其中\textbf{封闭孔}的体积。

\textbf{三种密度。}以不同的体积除质量即得三种密度：
\begin{equation}\label{eq:ch01c-rhoskel}
\rho_{\mathrm{skel}}=\frac{m}{V_{\mathrm{skel}}}
\end{equation}
\begin{equation}\label{eq:ch01c-rhopart}
\rho_{\mathrm{part}}=\frac{m}{V_{\mathrm{pore}}+V_{\mathrm{skel}}}
\end{equation}
\begin{equation}\label{eq:ch01c-rhobulk}
\rho_{\mathrm{bulk}}=\frac{m}{V_{\mathrm{app}}}
\end{equation}
既然分母满足 $V_{\mathrm{skel}}<V_{\mathrm{pore}}+V_{\mathrm{skel}}<V_{\mathrm{app}}$，
必有
\begin{equation}\label{eq:ch01c-order}
\rho_{\mathrm{skel}}>\rho_{\mathrm{part}}>\rho_{\mathrm{bulk}}
\end{equation}
\textbf{这个不等式是一个免费的自检工具}：任何报告如果违反它，必有一项测错或命名用错。

\textbf{两个派生量。}由骨架密度与颗粒密度可求得催化剂颗粒的\textbf{孔隙率}
（porosity）$p$：
\begin{equation}\label{eq:ch01c-poros}
p=\frac{V_{\mathrm{pore}}}{V_{\mathrm{pore}}+V_{\mathrm{skel}}}
 =1-\frac{\rho_{\mathrm{part}}}{\rho_{\mathrm{skel}}}
\end{equation}
由颗粒密度与堆积密度可求得催化剂堆积的\textbf{空隙分数}（空隙率，void fraction）
$\varepsilon$：
\begin{equation}\label{eq:ch01c-voidfrac}
\varepsilon=1-\frac{\rho_{\mathrm{bulk}}}{\rho_{\mathrm{part}}}
\end{equation}
\textbf{两式的分母不同，不可混用}：$p$ 描述的是\textbf{一颗颗粒内部}的开孔占颗粒体积的
分数，是材料的孔性；$\varepsilon$ 描述的是\textbf{床层}里颗粒与颗粒之间的空隙占外观体积
的分数，是装填状态的性质。同一种催化剂，$p$ 由制备决定，$\varepsilon$ 由装填方式决定。
两个分数的乘积才是开孔占床层外观体积的总分数。

\textbf{命名混乱提示。}英国标准协会（BSI）把\textbf{包含封闭孔}的骨架体积称为表观颗粒
体积（apparent particle volume），对应的密度称为表观颗粒密度；把\textbf{排除封闭孔}的
骨架体积称为真体积（true volume）或绝对粉末体积，对应的密度称为真密度或绝对粉末密度。
但在文献中，真密度与骨架密度常被当作同义词混用。\textbf{读文献数值时，若作者没有说明
封闭孔怎么处理，这个数就有约几个百分点的系统不确定度。}

\textbf{理论密度。}$\rho$ 还有一个更根本的定义：理论密度是指在原子水平上密堆积的、不含
任何孔与缺陷的颗粒的总质量与颗粒体积总和之比。但实际的固体催化剂往往是无定形的多种物质
的不均匀混合体，因此\textbf{理论密度这一概念对固体催化剂而言并没有多少实际意义}——它
只在讨论已知结构的纯晶相（如分子筛骨架）时才有确定值。

\textbf{包裹体积。}ASTM 引入了包裹体积（envelope volume）和包裹密度的概念来表征单个颗粒
的体积和密度：包裹体积就是假想在颗粒表面紧密包覆一层膜、膜内包容的体积。显然包裹体积
包含了颗粒的骨架体积和内孔体积，但由于表面的粗糙性，这一假想膜随观测尺度的不同而变化，
膜内也包覆了颗粒外表面的一些粗糙空隙。虽然包裹体积这一概念更接近实际，但从理论上它
\textbf{大于}本处所述的「颗粒体积」。\src{S1 §1.2.4，书内 p.75--77 / PDF p.94--96}
\end{principle}

\begin{figure}[htbp]
\centering
\begin{tikzpicture}[x=1mm,y=1mm]
% 比例说明
\node[anchor=west,font=\fontsize{8}{9.6}\selectfont] at (0,4)
  {示意比例：骨架 52\%、内孔 26\%、颗粒间空隙 22\%（按体积）};
% 三根条：堆积体积 / 颗粒体积 / 骨架体积
\fill[black!12]  (13.5,0) rectangle (48.5,-6);
\fill[black!30]  (48.5,0) rectangle (66.5,-6);
\fill[black!48]  (66.5,0) rectangle (81.5,-6);
\draw[thick] (13.5,0) rectangle (81.5,-6);
\fill[black!12]  (13.5,-10) rectangle (48.5,-16);
\fill[black!30]  (48.5,-10) rectangle (66.5,-16);
\draw[thick] (13.5,-10) rectangle (66.5,-16);
\draw[dashed,accentsoft] (66.5,-10) rectangle (81.5,-16);
\fill[black!12]  (13.5,-20) rectangle (48.5,-26);
\draw[thick] (13.5,-20) rectangle (48.5,-26);
\draw[dashed,accentsoft] (48.5,-20) rectangle (81.5,-26);
% 段内标签
\node[font=\fontsize{8}{9.6}\selectfont] at (31,-3)   {骨架};
\node[font=\fontsize{8}{9.6}\selectfont] at (57.5,-3) {内孔};
\node[font=\fontsize{8}{9.6}\selectfont] at (74,-3)   {空隙};
% 左侧标签
\node[anchor=west,font=\fontsize{8}{9.6}\selectfont] at (0,-3)  {堆积体积};
\node[anchor=west,font=\fontsize{8}{9.6}\selectfont] at (0,-13) {颗粒体积};
\node[anchor=west,font=\fontsize{8}{9.6}\selectfont] at (0,-23) {骨架体积};
% 右侧标签
\node[anchor=west,font=\fontsize{8}{9.6}\selectfont,align=left,text width=44mm] at (83.5,-3)
  {外观体积 $\rightarrow$ 堆积密度（含颗粒间空隙与内孔）};
\node[anchor=west,font=\fontsize{8}{9.6}\selectfont,align=left,text width=44mm] at (83.5,-13)
  {含内孔 $\rightarrow$ 颗粒密度（压汞法）};
\node[anchor=west,font=\fontsize{8}{9.6}\selectfont,align=left,text width=44mm] at (83.5,-23)
  {只含骨架 $\rightarrow$ 骨架密度（真密度）};
% 竖直参考线
\draw[dashed,accentsoft] (48.5,1.5) -- (48.5,-27.5);
\draw[dashed,accentsoft] (66.5,1.5) -- (66.5,-27.5);
\end{tikzpicture}
\caption{三种密度各自的分母是哪一份体积。读者应看\textbf{三根条的长度}：堆积体积最长
（骨架＋内孔＋颗粒间空隙），颗粒体积少了颗粒间空隙，骨架体积只剩固体本身；竖直虚线标出
三者的分界。由此可直接读出三种密度的定义式（式~\eqref{eq:ch01c-rhoskel}—
式~\eqref{eq:ch01c-rhobulk}）：\textbf{分子相同，分母依次变小，所以数值依次变大}。
示意图中的体积比例（骨架 52\%、内孔 26\%、空隙 22\%）只是为了让三段都看得见，不代表某种
具体催化剂的实测值；实际催化剂的内孔与空隙比例随制备与装填方式变化很大。\schemamark
\figsrc{据 S1（上册）§1.2.4，书内 p.75--76 / PDF p.94--95 的体积分解关系绘制。}}
\label{fig:ch01c-volume}
\end{figure}

\subsubsection{骨架密度（真密度）的测定}

骨架体积是催化剂颗粒固体骨架及其中包含的封闭孔的体积（此处参考了 ASTM 的定义）。它的
测定基于\textbf{阿基米德原理}：将一定质量 $m$ 的多孔颗粒（粉末）浸入可润湿的液体或气体
介质中，当润湿性介质完全进入颗粒的间隙和颗粒内的开孔后，颗粒置换出的介质体积即为颗粒的
骨架体积 $V_{\mathrm{skel}}$，于是 $\rho_{\mathrm{skel}}=m/V_{\mathrm{skel}}$。

测定方法有氦比重计法和比重瓶法两种。

（1）\textbf{氦比重计法}（helium pycnometry）以氦气作置换介质。将颗粒样品置于样品池中，
充高纯氦气至一定压力，然后连通参考池使氦气膨胀。样品池和参考池的体积已知，记录温度和两个
过程的平衡压力，根据\textbf{质量守恒}确定样品的体积，计算出骨架密度。氦分子小、惰性且不
被吸附，能进入绝大多数开孔；对氦气能渗透进去的材料，也可改用氮气或 \ce{SF6}。许多商品化
的全自动骨架密度计测量原理即基于氦比重计法。

（2）\textbf{比重瓶法}（specific gravity bottle method）在 Gay-Lussac 比重瓶中放置适量
的颗粒样品，通过差重法称量样品的质量及其置换出水的质量，计算骨架密度。比重瓶法设备简单、
精确度高，但操作略显烦琐。除水以外，还可以采用其他与试样润湿但不反应的液体介质（如煤油、
苯等）。\textbf{所测样品必须严格干燥；实验时样品须完全润湿并注意除气}——残留在孔内的
气泡会使置换体积偏小，从而把骨架密度算大。

至于「封闭孔算不算进骨架体积」，请回看 §1.2.4 的命名提示：含封闭孔的叫骨架密度（BSI 的
表观颗粒密度），排除封闭孔的叫真密度（BSI 的绝对粉末密度）。两种口径的差别就是封闭孔的
体积分数，对高封闭孔含量的样品可以到百分之几。\src{S1 §1.2.4 2)，书内 p.75--76 /
PDF p.94--95}

\subsubsection{颗粒密度的测定}

颗粒体积由颗粒内的孔体积 $V_{\mathrm{pore}}$ 和颗粒骨架体积 $V_{\mathrm{skel}}$ 两部分
组成，$\rho_{\mathrm{part}}=m/(V_{\mathrm{pore}}+V_{\mathrm{skel}})$。常压下，
\textbf{汞只能充填颗粒之间的间隙而不能侵入颗粒内部孔}，故可用\textbf{压汞法}测得颗粒体积，
求算颗粒密度；压汞法测得的颗粒密度也称\textbf{汞置换密度}。这里必须强调一句使用方法上的
话：\textbf{该值的大小与压汞法使用的压力有关}——压力升高后汞开始侵入较大的孔，被计入
「颗粒体积」的孔就减少，密度随之变大。因此报出汞置换密度时必须同时报出所用压力（或压力
上限）。

商品化的\textbf{包裹密度分析仪}采用另一种置换介质：用高流动性、均匀刚性的固体微球包覆所测
颗粒样品的外表面，并填充颗粒间隙空间（类似于代替了压汞法的置换介质汞），由此测定包裹密度。
包裹体积包含了骨架体积与内孔体积，由于表面粗糙性还会包住一部分外表面的粗糙空隙，所以
包裹密度略小于本处定义的颗粒密度。

本处使用的「颗粒密度」概念等同于 BSI 使用的「有效颗粒密度」（effective particle density）；
BSI 也有类似于 ASTM 的包裹体积与包裹密度概念。\src{S1 §1.2.4 3)，书内 p.76--77 /
PDF p.95--96}

\subsubsection{堆积密度的测定}

堆积密度定义为质量除以颗粒的堆积体积，堆积体积是一种\textbf{外观体积}（apparent volume），
$\rho_{\mathrm{bulk}}=m/V_{\mathrm{app}}$，因此堆积密度又称外观密度。它与颗粒的形状、装填
的方式及容器的形状有关。采用标准化的装填和测定方法或者仪器可以得到重复性很好的颗粒堆积
密度。依据测定方法的不同，堆积密度也有多种分类，如\textbf{振实堆积密度}、\textbf{松装堆积
密度}、\textbf{振动堆积密度}（vibrational packing density，ASTM）和\textbf{体积密度}
（bulk density）。测堆积密度的标准方法有 GB/T 6286—1986《分子筛堆积密度测定方法》、
ASTM D4180—13、ASTM D6683—14 等。

同样的命名分歧也出现在这里：ASTM D3766—05（2013）也使用体积密度（bulk density）一词表示
堆积密度；BSI 使用表观粉末密度（apparent powder density）；tapped density 和
tap density 也各有自己的定义。商品化仪器测定堆积密度时并不使用完全相同的标准，
\textbf{使用数据时必须注意是哪一个标准、哪一种装填方式}。

在工业生产实践中，固定床反应器催化剂装填还有\textbf{袋式装填密度}（sock loading density）
和\textbf{密相装填密度}（dense loading density）之分，它们是特定催化剂装填方式下工业催化剂
在反应器内的堆积密度：

\begin{itemize}
\item \textbf{袋式装填法}是通过一个接在催化剂料斗出口的帆布筒，靠自然重力直接把催化剂倒入
反应器的简单装填方法。适用于催化剂易积炭、床层易堵塞的反应过程。缺点是催化剂装填得不是很
均匀，容易产生沟流和热点。
\item \textbf{密相装填}是利用专门的装填器，通过空气推进或动能推进的方式将催化剂条均匀地
水平分散在催化剂床层截面。密相装填有效地避免了因催化剂条架桥造成的床层空隙，能够强化传质，
降低返混与沟流。\textbf{密相装填密度一般比袋式装填密度大 10\%—15\%}。
\end{itemize}

这 10\%—15\% 不是测量误差，而是\textbf{同一个床层里真实存在的差别}：它同时改变床层的
持气量、压降与停留时间分布。在反应器设计中把实验室松装密度直接用作工业装填密度，是
装填量算错、热点位置预测错误的常见原因。\src{S1 §1.2.4 4)，书内 p.77--78 /
PDF p.96--97}

\begin{table}[htbp]
\centering
\caption{三种密度：定义、测定方法、典型值与常见误用。读者应重点看第二列与第四列：
第二列决定这个数\textbf{测的是哪一种体积}，第四列是把它用错时最常出现的方式。
\supp{「典型值」一列为数量级参考（以 γ-氧化铝基载体与条状催化剂为例），不是原书数据；
具体数值随制备与成型工艺变化很大。}\src{原书无此汇总表，据 S1 §1.2.4，书内 p.75—78 /
PDF p.94—97 归纳。表中「典型值」一列及「常见误用」列为改编本归纳与教学设定，非源书实验数据。}}
\label{tab:ch01c-density}
{\srcsize
\begin{tabularx}{\columnwidth}{@{}lXXXX@{}}
\toprule
名称（文献别名） & 计入的体积（分母） & 测定方法与依据标准 & 典型数量级 & 常见误用 \\
\midrule
骨架密度（真密度；BSI：表观颗粒密度／绝对粉末密度） &
固体骨架体积 $V_{\mathrm{skel}}$，\textbf{含封闭孔}；排除封闭孔的口径称真体积 &
氦比重计法（气体膨胀、质量守恒）；比重瓶法（阿基米德原理，介质为水、煤油、苯等）。样品须严格干燥、完全润湿并除气 &
约 $3.0$--$3.2$ $\mathrm{g\,cm^{-3}}$ &
把「含封闭孔」与「排除封闭孔」两种口径混为一谈；把 XRD 精修得到的晶胞密度（理论密度）当作骨架密度 \\
颗粒密度（汞置换密度；BSI：有效颗粒密度） &
$V_{\mathrm{pore}}+V_{\mathrm{skel}}$，即骨架＋颗粒内\textbf{开孔} &
压汞法：常压下汞只充填颗粒间空隙、不侵入内孔；包裹密度分析仪用刚性微球填充颗粒间隙 &
约 $1.3$--$1.6$ $\mathrm{g\,cm^{-3}}$ &
把颗粒密度当作真密度使用；引用汞置换密度而不注明压汞压力（压力越高数值越大） \\
堆积密度（外观密度；体积密度；表观粉末密度） &
$V_{\mathrm{app}}=V_{\mathrm{void}}+V_{\mathrm{pore}}+V_{\mathrm{skel}}$，\textbf{含颗粒间空隙} &
标准量筒自由沉降（松装）、振实（GB/T 6286）、振动（ASTM D4180、D6683）；工业上另有袋式与密相装填密度 &
约 $0.6$--$0.8$ $\mathrm{g\,cm^{-3}}$ &
把堆密度当真密度；不写装填方式与依据标准——袋式与密相装填可相差 10\%—15\% \\
理论密度（晶胞密度） &
原子水平上密堆积、无孔无缺陷的固体体积 &
由晶体结构与晶胞参数计算，只对已知结构的纯晶相有确定值 &
随晶相而定 &
对无定形多组分催化剂使用这一概念；把它与骨架密度互换使用 \\
\bottomrule
\end{tabularx}}
\end{table}

\begin{measure}
\textbf{三个数各自的条件。}（1）\textbf{骨架密度}：样品必须严格干燥（表面水与孔内水都会
使质量偏大）、完全润湿并除气；气体置换法要记录两个平衡压力与温度；液体置换法要选与试样
润湿但不反应的介质。（2）\textbf{颗粒密度}：决定性的条件是\textbf{压力}——常压下汞只填
间隙，一旦加压到汞开始侵入孔，测到的就不再是颗粒体积；包裹密度法要说明所用微球的粒径与
流动性。（3）\textbf{堆积密度}：决定性的条件是\textbf{装填方式}（自由沉降、振实、振动、
工业袋式／密相）与容器形状；不同标准之间不可直接比较数值。
（4）\textbf{取样}：固定床催化剂的上下段、不同批次、运输前后，其堆密度与颗粒密度都可能有
系统差别，必须写明取样位置与是否经过筛分除去碎粉。\src{S1 §1.2.4，书内 p.75--78 /
PDF p.94--97}
\end{measure}

\begin{dataflow}
\textbf{从称重与压力数据到孔性判断。}（1）原始量：质量 $m$、置换出的介质体积（或氦气膨胀
前后的压力与池体积）、汞压入量与压力、量筒中的堆积体积。（2）按式~\eqref{eq:ch01c-rhoskel}
—式~\eqref{eq:ch01c-rhobulk} 得到三个密度。（3）用式~\eqref{eq:ch01c-order} 检查次序：
$\rho_{\mathrm{skel}}>\rho_{\mathrm{part}}>\rho_{\mathrm{bulk}}$ 必须成立，不成立就先
查命名与单位。（4）用式~\eqref{eq:ch01c-poros} 与式~\eqref{eq:ch01c-voidfrac} 分别算出
颗粒孔隙率 $p$ 与床层空隙率 $\varepsilon$。（5）判断：若 $p$ 明显小于由氮吸附测得的
比孔容所对应的孔隙率，说明颗粒密度测高了（多半是汞压过高，把一部分内孔算进了骨架）；
若 $\varepsilon$ 异常小，说明装填方式或碎粉含量有问题。
\src{S1 §1.2.4，书内 p.75--78 / PDF p.94--97}
\end{dataflow}

\begin{pitfall}
\textbf{误判一：把堆密度当真密度。}这是本节最高频的错误。堆密度的分母里含颗粒间空隙，
它的数值比骨架密度小好几倍；把它代入孔容、孔隙率或传质计算，会连带把后面所有的数算错。
判据：任何密度数值都必须能回答「分母里有没有颗粒间空隙」。

\textbf{误判二：把压汞法的颗粒密度当真密度。}汞置换密度含颗粒内开孔，它等于 BSI 的有效
颗粒密度，不等于骨架密度；而且它随压汞压力变化。

\textbf{误判三：不写方法与口径就引用数值。}同一名称在不同标准下定义不同（骨架密度 vs 真
密度、堆积密度 vs 体积密度 vs 表观粉末密度），文献值相差几个百分点到十几个百分点时，往往
不是样品差别而是命名差别。

\textbf{误判四：取样不代表整体。}催化剂床层上段与下段经历的温度、积炭与磨损不同，堆密度、
颗粒密度与强度都可能不同；用床层某一段的样品代表整个反应器，是工业测定的常见错误。此外，
样品若在取样时产生碎粉而未过筛，堆密度会偏大（细粉填隙）。\src{S1 §1.2.4，书内
p.75--78 / PDF p.94--97}
\end{pitfall}

\begin{crosscheck}
\textbf{四种互查，全部只用已经测过的数。}（1）\textbf{次序检查}：由
式~\eqref{eq:ch01c-order}，三个密度必须依次减小；违反次序说明命名或测量有误。
（2）\textbf{比孔容检查}：由三种密度可直接算出\textbf{比孔容}
\begin{equation}\label{eq:ch01c-vp}
v_{\mathrm{p}}=\frac{1}{\rho_{\mathrm{part}}}-\frac{1}{\rho_{\mathrm{skel}}}
\end{equation}
它应当与氮吸附法（或压汞法）测得的比孔容在允差内一致。不一致时按下面的顺序排查：先看
$v_{\mathrm{p}}$ 偏小还是偏大——偏小通常是颗粒密度测大了（压汞压力过高、或颗粒密度里
混入了封闭孔），偏大通常是骨架密度测小了（样品未完全干燥或除气不净）。
（3）\textbf{孔隙率与空隙率的量纲检查}：$0<p<1$、$0<\varepsilon<1$ 且两者互不包含；
把 $p$ 与 $\varepsilon$ 相加当作总孔率是错的。
（4）\textbf{与强度数据回查}：密相装填密度高的床层，其整体堆积抗压碎强度通常也更高
（架桥少、接触点多），若两者趋势相反，应检查装填方式或颗粒形状是否与报告一致。
\src{S1 §1.2.4，书内 p.75--78 / PDF p.94--97}
\end{crosscheck}

\subsection{催化剂机械强度的测定}

\begin{probout}
\textbf{本节回答}：催化剂的「强度」到底指什么；单颗粒强度与整体（堆积）强度为什么不能
互相替代；同一个样品为什么换一种加载方式就换一个数值；以及强度数据必须连同哪些条件报告。

\textbf{输出}：静态负荷下的\textbf{抗压碎强度}（单颗粒径向抗压碎力、点抗压碎力、整体堆积
抗压碎强度）与动态负荷下的\textbf{磨损性能}（磨耗率、磨损指数）。

\textbf{本节核心判断}：催化剂的机械强度是指其抵抗各种外界机械应力作用的能力。工业催化剂
从其生产、运输、装填直至工艺运转要经历多种应力，机械强度是判断其可靠性的重要物性技术指标
之一。为保证反应的顺利进行，主要关心的还是它\textbf{在实际工艺运转条件下}承受机械应力、
磨损和碰撞的能力。固体催化剂的工业应用主要有固定床和流化床两种形式，因而机械强度主要用
\textbf{静态负荷}和\textbf{动态负荷}测量：静态时的机械强度主要以\textbf{抗压碎强度}
表征，动态时则以\textbf{磨损性能}表征。有关测试方法大多实现了标准化，其设计原则都是尽可能
模拟催化剂在反应器内的实际工况，以保证强度数据对催化剂的应用有充分敏感性——\textbf{因此
强度数值只有在「加载方式＋试样形状＋判据」三者都相同时才可以比较}。
\end{probout}

\subsubsection{静态负荷：抗压碎强度}

\paragraph{单颗粒抗压碎强度}代表性标准有 HG/T 2782—2011《化肥催化剂颗粒抗压碎力的测定》、
HG/T 2783—1996（2010）《分子筛抗压碎力试验方法》、ASTM D4179 和 ASTM D6175—03（2013）。
单颗粒抗压碎强度测定使用颗粒强度测定仪，对处理后的催化剂样品颗粒\textbf{径向}（条形、
圆柱形样品）或者\textbf{点}（球形样品）接触施加压力，记录该颗粒被压碎时所施的力：

\begin{itemize}
\item \textbf{径向抗压碎力}：以颗粒径向方向的一个长度单位上所受的力表示，单位
$\mathrm{N\,cm^{-1}}$；
\item \textbf{点抗压碎力}：以颗粒直径方向上受的力表示，单位 N。
\end{itemize}

\textbf{单位不同就意味着物理意义不同}：径向值除以颗粒长度，反映的是条状催化剂侧壁的抗
压能力；点值不除以长度，反映的是球形颗粒整体的抗压能力。两者不能互相换算。为保证数据的
可靠性和代表性，需要测量\textbf{足够数量}的颗粒，数据以抗压碎力\textbf{均值和变异系数}
评价试样的抗压碎力及其均匀性——均值说明强度水平，变异系数说明成型质量的一致性；一批
均值合格但变异系数很大的催化剂，在床层中总有少数先碎的颗粒，而这些碎粉正是压降上升与
沟流的起点。\src{S1 §1.2.5 1)1)，书内 p.78 / PDF p.97}

\paragraph{整体堆积抗压碎强度}对固定床来讲，单颗粒强度并不能直接反映催化剂在床层中整体
破碎的情况，因而需要一种接近固定床真实情况的强度测试方法来表征催化剂的整体强度性能，即
\textbf{整体堆积抗压碎强度}；很多不规则形状的催化剂，其强度测试也只能采用这种方法。
代表性标准是 SHELL method SMS-1471，其基本方法如下：使用堆积抗压碎强度仪或自制设备，
测试的催化剂颗粒直径应不大于 6 mm；称量 $\qty{20}{cm^3}$ 处理后的催化剂样品放入一个截面
积为 $\qty{600}{mm^2}$ 的圆柱体不锈钢样品测量室；催化剂上覆盖直径 3—6 mm（因催化剂颗粒
尺寸而异）的\textbf{钢球} $\qty{5}{cm^3}$；不断增加的外力通过活塞施加到催化剂上，筛选并
称重在不同压力下压碎产生的\textbf{小于
420 µm 的粉末}（通过 ASTM 40 目标准筛）。堆积抗压强度 $P$ 定义为获得质量分数 0.5\% 的
420 µm 粉末时的压强（MPa）：

\begin{equation}\label{eq:ch01c-crush}
P=\frac{F}{A}
\end{equation}

式中 $F$ 为产生 0.5\% 粉末时施加在催化剂床上的力（N）；$A$ 为样品测量室的截面积
（$\mathrm{mm^2}$）。也可以通过绘制曲线（压碎粉末质量对压强）报告催化剂的堆积抗压强度。
此处已对照原页面图像（书内 p.79 / PDF p.98）：原书印「测试的催化剂颗粒直径应不大于
$6$ mm；称量 $20$ cm$^{3}$ 处理后的催化剂样品放入一个截面积为 $600$ mm$^{2}$ 的圆柱体
不锈钢样品测量室，催化剂上覆盖直径 $3\sim6$ mm（因催化剂颗粒的尺寸而异）的钢球
$5$ cm$^{3}$」，并明确 $P$ 以 MPa 计。原 PDF 文本层把「钢球」一类字样损伤为「钢球ｓｃ
耐」，本段据页面图像补全了钢球的直径与体积，并把此前误写的「钢制压板」改回原书的
\textbf{钢球}。\src{S1 §1.2.5 1)2)，书内 p.78--79 / PDF p.97--98}

\subsubsection{动态负荷：磨损性能}

\paragraph{磨耗率测定}磨耗率表征动态负荷条件下催化剂颗粒的抗磨损强度，代表性标准有
HG/T 2976—2011《化肥催化剂磨耗率的测定》、HG/T 3590—1999（2010）《制冷系统用分子筛
干燥剂抗磨耗性能的试验方法》及 ASTM D4058-96（2015）。其基本原理是：将经过预处理的催化剂
样品置于磨耗仪的磨耗筒（rotating drum）内，转动规定的时间，由于摩擦和磨蚀样品产生细小
碎颗粒（小于 850 µm，即通过 20 目标准筛），通过筛分称量计算磨耗率：

\begin{equation}\label{eq:ch01c-attrition}
A=\frac{m_{0}-m_{1}}{m_{0}}\times 100\%
\end{equation}

式中 $A$ 为磨耗率；$m_0$ 为磨耗前样品的质量；$m_1$ 为磨耗后样品的质量（以磨耗前后样品的
质量差为分子）。该法适用于片状、球状、挤压成型和不规则颗粒状催化剂。另外还有一种旋转管式
磨耗测试仪（rotating tube，Spence method）。

\paragraph{磨损指数测定}流化床工艺中的催化剂微球因高速气流的冲击而发生相互摩擦，磨损产生
的细粉会被气流夹带而损失。流化床催化剂的耐磨性能用\textbf{磨耗指数}（attrition index）
来表征，其测定依据 ASTM D5757—11（2017）。测定方法也称\textbf{高速空气喷射法}，适于测定
10—180 µm 的球体和不规则颗粒：将一定量的微球催化剂放入流化床磨损测定仪，利用特定流速的
喷射气流通过样品而发生磨耗，一定时间后产生的细粉（小于 20 µm 的颗粒）的质量分数即为
\textbf{空气喷射磨损指数}（air jet attrition index，AJI）。

\textbf{两个动态指标互不替代}：磨耗率用转筒的翻滚与滑动模拟固定床装卸与运转中的摩擦，
以 850 µm 为细粉界限；磨损指数用高速气流模拟流化床中的颗粒—颗粒、颗粒—器壁碰撞，以
20 µm 为细粉界限。同一个样品在两套方法下的排序可以不一致，因为它们模拟的是不同的磨损
机制（前者以表面磨损为主，后者以碎裂为主）。\src{S1 §1.2.5 2)，书内 p.79--80 /
PDF p.98--99}

\begin{table}[htbp]
\centering
\caption{四类强度测定的加载方式、结果表示与代表标准。读者应重点看第二列：
\textbf{试样形状与加载方式写在方法名里}，因而换一种加载方式就换一个数值，换一种结果表示
就不能互比。\src{原书无此类汇总表，据 S1 §1.2.5，书内 p.78—80 / PDF p.97—99 归纳。}}
\label{tab:ch01c-strength}
{\srcsize
\begin{tabularx}{\columnwidth}{@{}lXXX@{}}
\toprule
方法 & 加载方式与试样 & 结果表示 & 代表标准 \\
\midrule
单颗粒抗压碎强度 & 单颗条形／圆柱形颗粒\textbf{径向}对压，或球形颗粒\textbf{点}接触，加压至破碎 & 径向抗压碎力 $\mathrm{N\,cm^{-1}}$；点抗压碎力 N；报均值与变异系数 & HG/T 2782；HG/T 2783；ASTM D4179；ASTM D6175 \\
整体堆积抗压碎强度 & $\qty{20}{cm^3}$ 样品装入截面积 $\qty{600}{mm^2}$ 的样品室，上覆直径 $\qty{3}{mm}$--$\qty{6}{mm}$ 的钢球 $\qty{5}{cm^3}$，通过活塞逐级加压，筛出并称量 $<420$~µm 粉末 & 产生 0.5\% 质量分数粉末时的压强 $P=F/A$，MPa；或绘「压碎粉末质量—压强」曲线 & SHELL SMS-1471 \\
磨耗率（转筒法） & 预处理样品在磨耗筒中转动规定时间，由摩擦与磨蚀产生细粉 & 磨耗率 $A=(m_0-m_1)/m_0\times100\%$，细粉界限 $<850$~µm（20 目） & HG/T 2976；ASTM D4058 \\
磨损指数（空气喷射法） & 微球样品在流化床磨损仪中受特定流速喷射气流的冲击 & $\qty{5}{h}$ 试验后 $<20$~µm 细粉的质量分数（空气喷射磨损指数 AJI） & ASTM D5757 \\
\bottomrule
\end{tabularx}}
\end{table}

\emph{核对说明}：已对照原页面图像（书内 p.78—80 / PDF p.97—99，上册 §1.2.5）。四行与原文基本逐字相符，
代表标准 HG/T 2782、HG/T 2783、ASTM D4179、ASTM D6175、SHELL SMS-1471、
HG/T 2976、ASTM D4058、ASTM D5757 均见于原文。核对中改正两处：堆积强度行原写「上覆压板」，原书 p.79 为在催化剂上覆盖直径 $\qty{3}{mm}$--$\qty{6}{mm}$ 的钢球 $\qty{5}{cm^3}$；
磨损指数行原写「一定时间后」，原书 p.80 明确为 $\qty{5}{h}$ 试验。原书同节另有旋转管式磨耗测试仪（可同时测
4 个样品）与制冷系统用分子筛干燥剂标准 HG/T 3590，本表未收。


\begin{measure}
\textbf{强度数据的报告清单。}（1）\textbf{试样}：颗粒形状（球、条、片、三叶草、不规则）、
尺寸（直径与长度）、是否经过预处理（干燥、焙烧、筛除碎粉）、样品数。（2）\textbf{加载
方式}：径向还是点接触、加载速率、是否逐级加压、压板／转筒／喷射气流的规格与流速。
（3）\textbf{判据}：破碎的判定（力值突降、出现规定尺寸细粉的质量分数）与细粉界限
（420 µm、850 µm、20 µm 各不相同）。（4）\textbf{统计量}：均值＋变异系数（单颗粒）、
或压强值（堆积）。（5）\textbf{取样位置}：床层上段／中段／下段，运输前后。
缺少任何一项，强度数值都无法与其他实验室的数据比较。\src{S1 §1.2.5，书内 p.78--80 /
PDF p.97--99}
\end{measure}

\begin{dataflow}
\textbf{从力—位移曲线到一条可用的强度判据。}（1）原始量：颗粒强度仪给出压碎力（力—位移
曲线），堆积强度仪给出压强与各级细粉质量，磨耗仪给出磨耗前后的质量。（2）单颗粒：读每条
曲线的峰值（破碎力），按径向或点接触的定义换算成 $\mathrm{N\,cm^{-1}}$ 或 N，对几十颗
样品求均值与变异系数。（3）堆积：把「压碎粉末质量分数—压强」曲线读到 0.5\% 处，用
式~\eqref{eq:ch01c-crush} 换算成 MPa；曲线的\textbf{形状}本身也是信息——拐点早、斜率平
说明颗粒在低压下就大量破碎。（4）动态：用式~\eqref{eq:ch01c-attrition} 算磨耗率，或直接
读磨损指数。（5）判断：单颗粒均值高但变异系数大，与堆积强度偏低同时出现，通常意味着成型
工艺不均匀（条内密度梯度、黏结剂分布不匀），而不是原料强度不够。
\src{S1 §1.2.5，书内 p.78--80 / PDF p.97--99}
\end{dataflow}

\begin{pitfall}
\textbf{误判一：强度数据不写加载方式。}径向抗压碎力（$\mathrm{N\,cm^{-1}}$）与点抗压碎力
（N）单位不同、物理意义不同；同一个样品在两种加载方式下的数值不能换算，也不能比较大小。
引用文献强度值时，若不核对加载方式，很容易得出「国外催化剂强度高一倍」的错误结论。

\textbf{误判二：用单颗粒强度代表床层整体。}单颗粒强度不能直接反映催化剂在床层中整体破碎
的情况，这正是需要整体堆积抗压强度的原因；而不规则形状的催化剂甚至只能用堆积法。

\textbf{误判三：把磨耗率与磨损指数混为一谈。}两者模拟的是不同磨损机制、细粉界限也不同
（850 µm 与 20 µm），分别对应固定床与流化床，排序可以不一致。

\textbf{误判四：只看均值不看变异系数。}均值反映强度水平，变异系数反映均匀性；床层的
压降上升往往由少数先碎的颗粒触发。

\textbf{误判五：取样不代表整体。}催化剂床层上段与下段经历的应力、磨损和积炭不同，运输与
装填过程本身就会产生细粉；不写明取样位置的强度数据只代表来样，不代表反应器里的催化剂。
\src{S1 §1.2.5，书内 p.78--80 / PDF p.97--99}
\end{pitfall}

\begin{crosscheck}
\textbf{三类强度数据的互相回查。}（1）\textbf{单颗粒与堆积互查}：整体堆积抗压碎强度低而
单颗粒径向抗压碎力正常，通常不是颗粒本身弱，而是\textbf{床层}的问题——颗粒形状不规整、
装填不均、架桥导致应力集中。（2）\textbf{静态与动态互查}：抗压碎强度高但磨耗率也高的
样品是存在的（硬而脆）；两者同时偏低则提示成型密度不足或黏结剂失效。若这两项与堆积密度
的趋势不一致（密度高而强度低），应检查是否在成型时加入了过多润滑剂或造孔剂。
（3）\textbf{与工艺表现互查}：床层压降上升速度、细粉带出量与磨损指数、磨耗率应当同向；
若现场表现为「压降快速上升」而实验室磨耗率很低，应反过来怀疑取样位置（现场碎粉来自床层
上部或运输环节）或热态下的强度衰减（水热老化后的强度重测）。
\src{S1 §1.2.5，书内 p.78--80 / PDF p.97--99}
\end{crosscheck}

\srcnote{S1（上册）§1.2.2 3)（续）—§1.2.5，书内 p.61--80 / PDF p.80--99。本节涉及的
ASTM／HG 标准编号与仪器量程取自原书表 1--10 与正文；原 PDF 为 OCR 文本层，公式与表格数字
均按 OCR 上下文与学科标准形式重建，存疑处已在正文用「待核对」标出。}
